\begin{definition}[Код Прюфера]
Вершины дерева помечены числами \(1,\ldots,n\), \(n\ge2\).
Пока остаётся больше двух вершин, удаляем лист с наименьшей меткой
и записываем метку его единственного соседа. Полученная последовательность
длины \(n-2\) называется кодом Прюфера.
\end{definition}
\begin{lemma}[Частоты меток]
В коде дерева метка \(v\) встречается ровно \(d_T(v)-1\) раз.
В каждом оставшемся дереве листья — ровно те его вершины,
которые отсутствуют в ещё не обработанном суффиксе кода.
\end{lemma}
\begin{proof}
Рёбра у \(v\) удаляются по одному. Пока \(v\) не удалена,
удаление соседнего листа записывает \(v\) и уменьшает её степень на 1.
Последнее инцидентное ребро исчезает при удалении самой \(v\)
или остаётся последним ребром дерева; в обоих случаях метка \(v\)
за него не записывается. Итого записей \(d_T(v)-1\).
То же рассуждение применимо к любому промежуточному дереву;
степень 1 равносильна нулю будущих записей.
\end{proof}
\textbf{Декодирование.} Для произвольного слова \(a_1,\ldots,a_{n-2}\)
поддерживаем множество не удалённых меток \(U\). На шаге \(t\)
выбираем наименьшую метку \(x\in U\), отсутствующую среди
\(a_t,\ldots,a_{n-2}\), проводим ребро \(xa_t\), удаляем \(x\) из \(U\)
и переходим к следующей позиции. Последние две вершины соединяем ребром.
\begin{lemma}[Корректность декодирования]
Алгоритм всегда выполним и строит дерево с заданным кодом.
\end{lemma}
\begin{proof}
Перед шагом \(t\) осталось \(n-t+1\) меток и \(n-t-1\) позиций кода.
Все метки суффикса ещё находятся в \(U\), поскольку ранее удалялись
только отсутствовавшие в более длинном суффиксе. Следовательно,
не менее двух меток \(U\) в суффиксе отсутствуют и нужная \(x\)
существует. При этом \(x\ne a_t\).

Если рассматривать шаги в обратном порядке, начинаем с последнего ребра
и каждый раз добавляем новую вершину одним ребром к уже построенному
графу: \(a_t\) в момент удаления \(x\) ещё не удалена.
Так сохраняются связность и отсутствие циклов, значит, построено дерево.
В оставшейся части этого дерева каждая вершина имеет степень
1 плюс число её вхождений в оставшемся коде: каждое такое вхождение
прикрепляет к ней один удаляемый ранее лист, а одно ребро связывает её
с более поздней вершиной либо является последним ребром.
Следовательно, на каждом шаге выбранная \(x\) действительно есть
наименьший лист оставшегося дерева. Повторное кодирование даёт \(a_t\)
на каждом шаге, то есть исходное слово.
\end{proof}
\begin{theorem}[Биекция Прюфера и формула Кэли]
Кодирование задаёт биекцию между деревьями на \([n]\) и словами
длины \(n-2\) над алфавитом \([n]\). Поэтому число помеченных деревьев
при \(n\ge2\) равно \(n^{n-2}\).
\end{theorem}
\begin{proof}
Декодирование любого слова даёт дерево с этим кодом.
Если начать с дерева, лемма о частотах показывает, что декодер
на каждом шаге выбирает именно лист, удалённый кодером, и возвращает
все исходные рёбра. Таким образом, алгоритмы взаимно обратны.
У каждой из \(n-2\) позиций \(n\) независимых вариантов, всего \(n^{n-2}\).
\end{proof}
При \(n=2\) код пуст и дерево одно; при \(n=1\) дерево также одно,
этот случай рассматривается отдельно.
\begin{example}
Для звезды с центром 4 на \(\{1,2,3,4,5\}\) код равен \((4,4,4)\).
Декодер последовательно присоединяет к 4 вершины \(1,2,3\), затем
соединяет последние \(4,5\). Частота 4 равна \(3=d(4)-1\).
\end{example}
